\section{工具的对照与校准}\label{sec:comparison}

\subsection{对照表}

\begin{table}[H]
\centering
\caption{两类工具的对照}
\label{tab:contrast}
\small
\begin{tabular}{L{3.2cm}L{5.5cm}L{5.5cm}}
\toprule
 & \textbf{消费税（鞭炮、焰火）} & \textbf{禁放／限放} \\
\midrule
设立理由 & 1994 年即含“危害……生态环境”，写法“寓禁于征” & 改善大气环境质量（东营办法第一条首列） \\
设定层级 & \textbf{国家}（国务院令） & \textbf{市级}政府规章或通告 \\
城市维度分辨率 & \textbf{零}（全国统一税率） & \textbf{高}（逐城划定） \\
时间维度分辨率 & \textbf{零}（全年同率，且在生产环节征收） & \textbf{高}（逐时段，如“除夕至初五、正月十五”） \\
最后一次校准 & \textbf{三十余年未调整} & \textbf{逐城逐年调整}（东营：2021 暂行办法 $\rightarrow$ 2024 办法） \\
成本形态 & 征管成本极低；负担以涨价形式\textbf{可见地}落在消费者 & 巡查、蹲守、查处；成本是\textbf{警力与执法工时}，不进入价格、也不单独核算 \\
制度位置 & 挂在全国消费税征管上 & 写进以生态环境为主的多部门大气方案、与量化工业减排同文考核；牵头单位公安、应急管理、生态环境等均有，相当一部分文件未写明 \\
\bottomrule
\end{tabular}
\end{table}

\textbf{这张表的最后三行是本文的论点。} 两类工具都有环境理由；区别在于一个在结构上无法对准损害的两个维度，另一个可以；而可以的那个，其成本不进价格、也不被单独核算。

\subsection{庇古校准：算不出来，以及为什么算不出来}\label{sec:pigou}

一个自然的追问是：现行税负相当于边际损害的几分之几？

以“每公斤烟花”为共同分母，算式为
\begin{equation}
\text{税} = t\,P, \qquad \text{基准} = \mathrm{EF}\times \mathrm{NEC\_SHARE}\times C,
\label{eq:calib}
\end{equation}
其中 $P$ 为出厂价；EF 为排放因子（emission factor，kg \pmm{2.5} / kg 净火药量）；\texttt{NEC\_SHARE} 为净火药（net explosive content，NEC）占商品重量的比例；$C$ 为每公斤 \pmm{2.5} 的价值；VSL 为统计生命价值（value of a statistical life）。参数取值见附录~\ref{app:calibration}。

计算结果：

\begin{table}[H]
\centering
\caption{两条校准路线的结果}
\label{tab:calib}
\small
\begin{tabular}{L{3.0cm}R{1.6cm}R{2.3cm}R{2.3cm}R{1.8cm}R{2.0cm}}
\toprule
\textbf{路线} & \textbf{组合数} & \textbf{比值区间} & \textbf{中位} & \textbf{税$\geq$基准} \\
\midrule
A 避免治理成本 & 1,296 & 0.044—57.1 & 0.254 & 19\% \\
B 健康损害 & 4,096 & 0.0020—6.59 & 0.0225 & 11.7\% \\
\bottomrule
\end{tabular}
\end{table}

\textbf{两条路线都符号不稳，比值跨约三个数量级。} 中位数低于 1（路线 A 为 0.174、B 为 0.0154），但两条路线的上尾都越过 1（A 有 17\%、B 有 6.8\% 的组合给出“税 $\geq$ 基准”）。因此既不能说“税负一定偏低”，也不能报单点——\textbf{中位数给方向，区间否定单点，两者必须同时出现。}

这里必须交代三个陷阱，它们比结果本身更重要。

\textbf{陷阱一：路线 A 概念不适用。} “单位治理成本”预设排放可就地捕集（点源：烟囱、除尘器）。烟花 \pmm{2.5} 是弥散释放，无从捕集，唯一的减排就是不烧。因此 A 不是“更保守的估计”，而是\textbf{概念上用错了地方}。

\textbf{陷阱二：单位不匹配。} EF 的分母是净火药量，而 $P$ 是每公斤商品价格。两者之间必须有一个换算系数 \texttt{NEC\_SHARE}，而它没有来源。修正这一处之后，路线 B 的结论从“税始终低于基准”翻转为“11.7\% 的组合下税高于基准”。\textbf{修一处错误就翻转结论，说明结论对参数假设过度敏感。}

\textbf{陷阱三：短期暴露不能套长期系数。} 若把一天 $+43$ \ugm 当作年均升高 43 \ugm{} 代入长期剂量反应，会把损害放大数百倍。本文只用短期（每日死亡）系数，且只在那一天上计。

\textbf{这个失败本身有内容。} 主导不确定性的是扩散参数 $DK$（2e3—2e5，跨两个数量级，**自设**）；其余参数——EF、\texttt{NEC\_SHARE}（0.098—0.320，Tang et al., 2012 实测）、$P$、剂量反应、VSL——都有来源；而有可靠来源的四个参数（EF、$P$、剂量反应、VSL）恰恰不是主导项。更关键的是，\textbf{这两个参数之所以无法取值，原因正是它们必须按城市与时段索引}：同样的火药量，在静稳的盆地城市与在海滨城市的浓度贡献不同，在初一夜里与在平日也不同。

于是本文得到一个与第~\ref{sec:resolution}~节同向的推断：\textbf{对这个污染源，“每公斤的边际损害”不是常数，而是一个城市 $\times$ 时段上的函数。任何全国统一、按年适用的税率，在结构上都无法表达它。} 因此，“把这个税校准到边际损害”从来不是一个可选项；问题不在校准得不够好，而在“全国统一、按年适用”这个工具形式没有对应的边际损害参数。

\begin{quote}
这是一处推断，不是测量结果。本文能证明的是比值算不出来，以及算不出来的原因；“因此校准不可能”是由此推出的解释，不是被估计出的参数。
\end{quote}

\subsection{成本一侧：能说什么，不能说什么}\label{sec:cost}

执法数据的采集结果与预期相反：\texttt{data/enforcement.csv} 覆盖 21 个城市、70 行，其中 \texttt{cases} 只有 28 格有值、仅 8 格是“含违规燃放”的查处案件数；而 \textbf{\texttt{budget\_wan} 列在 21 个城市中全为空}——没有任何一个城市公开过禁放工作经费的年度数字。

能用的只有“专项”口径的人力投入：

\begin{table}[H]
\centering
\caption{禁放执法的人力投入（仅专项口径）}
\label{tab:enforcement}
\small
\begin{tabular}{L{2.2cm}L{4.2cm}L{3.6cm}R{1.6cm}C{1.4cm}}
\toprule
\textbf{城市} & \textbf{口径} & \textbf{人力投入} & \textbf{年份} & \textbf{档} \\
\midrule
武汉 & 禁鞭专项 & 5.6 万人次 & 2023 & N \\
郑州 & 禁限放专项 & 2.5 万人次 & 2026 & N \\
青岛 & 专项 & 2.7 万／0.5 万人次 & 2022／2023 & N \\
\midrule
石家庄 & \textbf{全市春节安保总量}（非专项） & 5.19 万人次 & 2022 & N（来源失效） \\
\bottomrule
\end{tabular}

\vspace{0.3em}
{\small 注：证据档 N 表示来源为媒体或官方通报、未经一手核实。石家庄一行\textbf{不能与前三行并读}：5.19 万是全市春节安保总量（含交通、治安、大型活动），不是禁放专项，列在此处是为了说明口径混用会把一个更大的数字误当成禁放成本。武汉 2023 年的 5.6 万含公安民警 3 万余人次，且同年是疫情后首个春节与全域禁放首年，案件数由 2021 年的 27 起跳至 1,187 起，多项因素叠加，不可单独解读。}
\end{table}

2011 年《检察日报》关于北京禁放的报道提供了一个更有说明力的细节：投入约 286 万人次管控力量，同时明确写着“直接的钱花得不多，主要是人力造成的间接成本”“工时的具体统计数字缺失”。\textbf{这是媒体来源，只能作为线索，不能作为数据。} 但它与本文的独立发现一致：\textbf{这项政策的主要成本是工时，而工时不被统计。}

因此本文给出的是\textbf{外部性代价与执法成本的下界}，不声称算出了最优政策强度。\textbf{成本不被核算这一事实本身，是本文的发现之一}——21 城无一公开经费、查处案件数只有个位数可用、人力口径混杂，这三件事合起来说明：这项政策的成本既不在预算中单列，也不在统计中留痕。
